% !PRINTED 506
\weberpage{506}

%<<CONT
darstellen, was in der That eine Gruppe ist, wie aus $\pi_1^{-1}\varkappa\pi_1\pi_1^{-1}\varkappa_1\pi_1 = \pi_1^{-1}\varkappa\varkappa_1\pi_1$ zu ersehen ist. Daraus folgt zugleich, dass $\pi_1^{-1}Q\pi_1$ mit $Q$ \gesperrt{isomorph} ist (\S.\,149).

Die Gruppen, zu denen die conjugirten Functionen (10) gehören, nämlich
\weq{12}{Q,\quad \pi_1^{-1}Q\pi_1,\quad \pi_2^{-1}Q\pi_2,\ldots \pi_{j-1}^{-1}Q\pi_{j-1}}
nennen wir \gesperrt{conjugirte Theiler} von $P$ oder kurz conjugirte Gruppen\footnote{Auch der Name \glqq{}gleichberechtigte Untergruppe\grqq{} ist dafür im Gebrauch.}.

Wenn $\pi$ eine beliebige Permutation aus $P$, also von der Form $\varkappa\pi_r$ ist, so ist die Gruppe $\pi^{-1}Q\pi$ immer unter den conjugirten Gruppen (12) enthalten, nämlich $= \pi_r^{-1}\varkappa^{-1}Q\varkappa\pi_r = \pi_r^{-1}Q\pi_r$.

Die Ableitung von $\pi^{-1}Q\pi$ aus $Q$ heisst auch eine \gesperrt{Transformation der Gruppe $Q$ durch $\pi$} und $\pi^{-1}Q\pi$ eine \gesperrt{aus $Q$ transformirte Gruppe}.

Zur Bildung der Permutationen der conjugirten Gruppen $\pi^{-1}Q\pi$ führt folgende einfache Regel:

6. \gesperrt{Man erhält die Permutation $\pi^{-1}\varkappa\pi$ dadurch, dass man in den Cyklen von $\varkappa$ die Permutation $\pi$ ausführt.}

Um die Regel zu beweisen, sei, in die Cyklen zerlegt (\S.\,153),
\[
\varkappa = (\alpha,\beta,\gamma\ldots)(\alpha',\beta',\gamma'\ldots)\ldots
\]
und es sei
\[
\pi =
\begin{pmatrix}
\alpha, & \beta, & \gamma & \ldots \alpha', & \beta', & \gamma' & \ldots \\
\alpha_\pi, & \beta_\pi, & \gamma_\pi & \ldots \alpha'_\pi, & \beta'_\pi, & \gamma'_\pi & \ldots
\end{pmatrix}.
\]
Durch $\pi^{-1}$ geht $\alpha_\pi$ in $\alpha$ über, durch $\varkappa$ geht $\alpha$ in $\beta$ über und durch $\pi$ wird $\beta$ in $\beta_\pi$ übergeführt. Durch $\pi^{-1}\varkappa\pi$ geht also $\alpha_\pi$ in $\beta_\pi$ über. Da dieselbe Betrachtung auf $\beta_\pi, \gamma_\pi\ldots$ u.\,s.\,f. anwendbar ist, so folgt:
\[
\pi^{-1}\varkappa\pi = (\alpha_\pi,\beta_\pi,\gamma_\pi\ldots)(\alpha'_\pi,\beta'_\pi,\gamma'_\pi\ldots),
\]
wodurch die Regel 6. erwiesen ist.

Von der Zerlegung in Nebengruppen können wir noch eine Anwendung machen auf den Beweis des Satzes:

7. \gesperrt{Der Grad einer transitiven Permutationsgruppe von $m$ Ziffern ist immer ein Vielfaches von $m$.}
